\chapter{SONUÇLAR VE ÖNERİLER}
\label{ch:sonuc}

Bu bölümde ara dönem itibarıyla çalışmanın genel değerlendirmesi yapılmakta, karşılaşılan zorluklar
anlatılmakta ve dönemin kalanı için iş planı sunulmaktadır.

\section{Genel Değerlendirme}

Ara rapor tarihinde, dönem başında hedeflenen sistem uçtan uca çalışmakta ve canlıda
bulunmaktadır: talepler etiketlenmekte, cevap resmî belgelerde aranmakta, kaynaklı bir taslak
yazılmakta ve gönderim kararı test edilmiş kodla verilmektedir. Öğrenci ve personel arayüzleri,
birime göre yetkilendirme, değiştirilemeyen talep geçmişi, istek sınırları ve saldırılara karşı
savunmalar tamamlanmıştır. Ayar için hiç kullanılmayan iki kilitli test setinde insana gitmesi
gereken taleplerden hiçbiri kaçırılmamış, doğru belge her seferinde ilk üç sonuç içinde bulunmuş ve
12 saldırı denemesinin tamamı engellenmiştir.

Çalışmanın en önemli tasarım sonucu, bir dil modelinin güvenilir bir sistemin parçası olabilmesi için
\textbf{kararın modelden ayrılması} gerektiğidir. Model etiketleri ve taslağı üretir; gönderim kararı,
kaynak kartları ve alıntı denetimi kodda yapılır. Bu sayede her kural ayrı ayrı test edilebilmekte,
her karar öğrenciye ve personele gerekçesiyle açıklanabilmektedir.

\section{Karşılaşılan Zorluklar}

\begin{itemize}
  \item \textbf{Ücretsiz katman sınırları:} İlk kullanılan model günde yalnızca 20 isteğe izin
        verdiği için değerlendirmeler yarım kalmış, daha yüksek kotalı bir modele geçilmiştir.
        Canlı sistemde modellerin sık sık meşgul olması, üç modelli dayanıklı zincirle çözülmüştür.
  \item \textbf{Eşik seçimi:} Benzerlik skorunun cevaplanabilir ve cevaplanamaz talepleri ayırmadığı
        görülmüş (Şekil~\ref{fig:skor}), karar tek bir eşik yerine birden çok kontrole dayandırılmıştır.
  \item \textbf{Bilgi tabanındaki boşluklar:} İlk kilitli test çalıştırmasında bir SOFRA şifresi
        sorusu yanlış belgeyle cevaplanmıştır. Resmî Bilişim Hesapları belgesi eklenerek ve ``acil talep
        her zaman insana gider'' kuralı yazılarak bu boşluk kapatılmıştır.
  \item \textbf{İki sorun içeren talepler:} Bir talepte iki ayrı sorun olduğunda biri atlanabilmektedir
        (Bölüm~\ref{ch:bulgular}). Bu sorun kendi testlerimde bulunmuş ve düzeltme sıraya alınmıştır.
  \item \textbf{Ücretsiz barındırma:} Sunucu 15 dakika ziyaretçi olmazsa uyumakta, ilk açılış
        yaklaşık 30 saniye sürmektedir.
\end{itemize}

\section{Dönemin Kalanı İçin İş Planı}
\label{sec:isplani}

Şekil~\ref{fig:gantt} tamamlanan işleri ve kalan dönemin planını haftalık olarak göstermektedir.
Planlanan her iyileştirme önce geliştirme setinde ölçülecek, kilitli setler yalnızca dönem sonunda
bir kez çalıştırılacaktır.

\begin{figure}[!ht]
\centering
\scalebox{0.86}{%
\begin{ganttchart}[
    hgrid, vgrid={*1{gray!25, dotted}},
    x unit=0.62cm, y unit chart=0.55cm, y unit title=0.55cm,
    title label font=\scriptsize\bfseries, title height=1,
    bar label font=\scriptsize, group label font=\scriptsize\bfseries,
    bar/.append style={fill=s1, draw=none}, bar height=0.55,
    group/.append style={fill=koyu}, group height=0.25,
    milestone/.append style={fill=bseu, draw=none}, milestone label font=\scriptsize\bfseries,
    today=9, today offset=0.5, today label=Ara rapor, today label font=\scriptsize\color{bseu},
    today rule/.style={draw=bseu, dashed, thick},
    bar label node/.append style={align=right}
  ]{1}{16}
  \gantttitle{Eyl}{1} \gantttitle{Ekim}{4} \gantttitle{Kasım}{5} \gantttitle{Aralık}{4} \gantttitle{Ocak}{2} \\
  \gantttitlelist{1,...,16}{1} \\
  \ganttgroup{Tamamlanan işler}{1}{9} \\
  \ganttbar{Konu seçimi, problem araştırması}{1}{2} \\
  \ganttbar{Etiket kural kitabı ve 80 talepli geliştirme seti}{2}{2} \\
  \ganttbar{Sınıflandırıcı, kilitli test seti, ilk ölçümler}{2}{3} \\
  \ganttbar{Belge arama (RAG), kaynaklı cevap, karar kuralı}{3}{3} \\
  \ganttbar{API, personel kuyruğu, PostgreSQL + pgvector}{3}{4} \\
  \ganttbar{CI, Docker, React arayüzü, personel girişi}{3}{5} \\
  \ganttbar{Canlıya alma, saldırı testleri, gerçekçi set}{5}{5} \\
  \ganttbar{Ön bilgi, genel tavsiye, model zinciri, geçmiş}{6}{8} \\
  \ganttmilestone{Ara rapor}{9} \\
  \ganttgroup[group/.append style={fill=s2}]{Planlanan işler}{9}{16} \\
  \ganttbar[bar/.append style={fill=s2}]{İki sorunlu talep düzeltmesi, tam değerlendirme}{9}{10} \\
  \ganttbar[bar/.append style={fill=s2}]{Karma arama (tam metin + vektör, RRF)}{10}{11} \\
  \ganttbar[bar/.append style={fill=s2}]{200+ talepli set, model karşılaştırması}{11}{13} \\
  \ganttbar[bar/.append style={fill=s2}]{İzleme (OpenTelemetry), maliyet paneli}{13}{14} \\
  \ganttbar[bar/.append style={fill=s2}]{Kilitli setlerde nihai ölçüm, final rapor, sunum}{15}{16}
\end{ganttchart}}
\caption{İş planı (1. hafta: 28 Eylül 2026); kesikli çizgi ara rapor tarihini gösterir}
\label{fig:gantt}
\end{figure}

\section{Gelecekteki Çalışmalar için Öneriler}

Bu çalışmanın ilerleyen aşamalarında;

\begin{itemize}
  \item İki sorun içeren talepler ayrı ayrı ele alınarak her sorun kendi birimine yönlendirilebilir.
  \item Anahtar kelime (PostgreSQL tam metin arama) ve vektör arama RRF ile birleştirilerek, hit@1 ve
        MRR ölçütlerindeki etkisi ölçülebilir \cite{cormack2009}.
  \item Değerlendirme seti 200'ün üzerine çıkarılarak modeller (yerel bir açık kaynak model dâhil)
        doğruluk, kaçırılan escalation, 1000 talep başına maliyet ve gecikme bakımından
        karşılaştırılabilir.
  \item Cevap kalitesi ikinci bir dil modeliyle puanlanabilir \cite{zheng2023}, sınıf arkadaşlarından
        toplanan gerçek talepler yeni bir görülmemiş test seti olarak kullanılabilir.
  \item İkinci bir etiketleyiciyle etiketler arası uyum ölçülebilir.
  \item Üniversite tek oturum açma (SSO) sistemi ve personel cevabının öğrenciye e-postayla
        bildirilmesi eklenerek sistem gerçek kullanıma hazırlanabilir.
\end{itemize}

Bu öneriler doğrultusunda yapılacak çalışmalar, sistemin hem daha çok öğrenciye hemen yardım
etmesine hem de güvenlik sonuçlarının daha büyük ve gerçekçi verilerle doğrulanmasına katkı
sağlayacaktır.
